
\section{Methods and Experimental Approach}
\label{section_methods}

Temporal distribution shift manifests differently across modalities, tasks, and time scales, yet existing empirical studies typically vary a single axis at a time, leaving open how \emph{architectural design}, \emph{label structure}, and \emph{drift mechanism} jointly shape temporal robustness (Section~\ref{section_background}). Our setup targets exactly this cross-cutting comparison: we combine three long-range, time-indexed datasets with a diverse suite of neural architectures spanning multiple inductive biases, all evaluated under the unified temporal protocol of Section~\ref{section:evaluation_framework} (temporal drift matrices), so that architectural differences, rather than differences in splitting or evaluation, drive the observed robustness patterns.

\subsection{Datasets}
Our empirical analysis spans three qualitatively distinct temporal scenarios, chosen to cover different modalities, tasks, and sources of distribution shift. Each dataset captures multiple decades of real-world temporal evolution, making it suitable for cross-temporal evaluation.

\subsubsection{Yearbook: A Century of Portraits}
The \texttt{Yearbook} dataset \iflncs\cite{ginosar2015century}\else\cite{yearbook_ginosar_2019}\fi{} contains 37{,}921 frontal portraits of American high-school seniors from \val{yearbook/cutoff-first} to \val{yearbook/cutoff-last}, which we align and process as 3-channel $32\times32$ tensors. Although acquisition is largely standardized, stylistic attributes (e.g., hairstyles, clothing, accessories) vary substantially across decades. The dataset provides binary sex labels that are approximately balanced over time. It is a canonical benchmark for temporal shift: Wild-Time \cite{yao2022wildtime} uses a pre-/post-1970 split and reports marked out-of-distribution degradation, and Modyn \cite{boether_modyn_2025} observes accuracy decay as the train--test time gap grows. We reproduce this trend (Fig.~\ref{fig:mean_matrices}), consistent with covariate and concept drift.

\subsubsection{Amazon Reviews 2023: E-commerce}
The \texttt{Amazon Reviews} dataset \cite{hou2024bridging} comprises 571.54 million reviews across 33 product categories, spanning May 1996 to September 2023. Each review has a timestamp, star rating, and free-text content. We cast a review-level sentiment regression task, predicting the 1--5 rating from the text, which exhibits strong covariate and concept drift due to evolving language, consumer behavior, and platform usage. We focus on seven categories, restrict to 2014--2023, and draw a stratified sample of 300{,}000 reviews.

\subsubsection{arXiv: Scientific Discourse}
The \texttt{arXiv} dataset \cite{arxiv2024} defines a multi-label task over 2{,}866{,}787 title--abstract records annotated with 176 subject categories. We concatenate titles and abstracts, keep seven leaf categories (\cref{app:arxiv}), and use \val{arxiv/cutoff-first}--\val{arxiv/cutoff-last} submissions whose categories fall within them (papers may carry several). Category mix and terminology shift, inducing label and covariate drift\iflncs\else\ \cite{holzinger2024analysis}\fi. Wild-Time \cite{yao2022wildtime} reports roughly 20\% degradation under temporal splits for a related arXiv task, and the gap is not substantially closed by domain generalization or continual learning methods, motivating our architectural comparison.

\subsection{Model Implementation}

We evaluate architectures spanning different inductive biases to understand how model design affects temporal robustness. Rather than covering the full architecture landscape, we select families that span the spectrum of inductive-bias strength: from simple baselines without structural assumptions that establish lower bounds on performance, to modern architectures that incorporate structural priors, to pre-trained models that leverage large-scale external data. This spread is what lets us attribute robustness differences to the priors themselves.

\subsubsection{Image Classification}

For the \texttt{Yearbook} dataset, we evaluate four model families trained on the data, each at three sizes (small, medium, large), for \val{yearbook/scratch-models} architectures, complemented by \val{yearbook/frozen-models} pretrained vision encoders used as frozen backbones.

The four families differ in the spatial prior they encode. At one extreme, the \emph{multilayer perceptron} (\texttt{MLP}) flattens the image and applies only fully connected layers, imposing no spatial structure. The \emph{convolutional network} (\texttt{CNN}) builds in locality and translation equivariance through convolutional filters with batch normalization and pooling, and the \emph{residual network} (\texttt{ResNet}) extends it with skip connections \cite{he2016resnet} that ease optimization at greater depth. At the other extreme, the \emph{vision Transformer} (\texttt{ViT}) drops convolution for self-attention over patches \cite{dosovitskiy2021vit}, a far weaker spatial prior, with a small patch size suited to the $32\times32$ inputs. Across all four, the small, medium, and large variants scale depth and width.

Finally, we assess \emph{transfer learning} using pretrained vision encoders trained on large-scale curated or web-scale datasets. The underlying hypothesis is that representations learned from diverse data may exhibit stronger temporal robustness than features learned solely from historical portraits. We consider \val{yearbook/frozen-models} pretrained models. The self-supervised \texttt{DINOv2-S} \cite{oquab2024dinov2} and \texttt{DINOv3-S} \cite{simeoni2025dinov3} are pretrained on large curated image collections; \texttt{CLIP-B32} \cite{radford2021clip} and \texttt{SigLIP-B} \cite{zhai2023siglip} use contrastive image-text pretraining, the latter with a sigmoid loss; \texttt{ConvNeXt-S} \cite{liu2022convnext}, \texttt{ResNet50-IN} \cite{he2016resnet}, and \texttt{ViT-S16-IN21k} \cite{dosovitskiy2021vit} are trained with supervised \texttt{ImageNet} labels, the last on \texttt{ImageNet-21k}; and \texttt{MAE-B} \cite{he2022mae} and \texttt{EVA02-B} \cite{fang2024eva02} are pretrained with masked image modeling. In all cases, we freeze the pretrained backbone and train only a linear classification head, isolating the contribution of pretrained representations from the effects of fine-tuning dynamics.




\subsubsection{Text Models}

For the \texttt{Amazon Reviews} and \texttt{arXiv} datasets, we evaluate model families that differ in their inductive biases for representing textual structure. The same architectures serve both datasets, differing only in output layer and loss, with \texttt{Amazon Reviews} using regression under a weighted mean squared error and arXiv multi-label classification under a weighted binary cross-entropy.

All four families operate on cached \texttt{RoBERTa} token embeddings. The \emph{feed-forward baseline} (\texttt{FFN}) averages them into a single 768-dimensional vector and passes it through a small MLP head that discards order entirely, with variants scaling the hidden width from 128 to 2048. The \emph{convolutional model} (\texttt{TextCNN}) \cite{kim2014convolutional} applies one-dimensional convolutions to capture local n-gram patterns, scaling the number and width of its filters. The \emph{recurrent} models read the sequence token by token, as a single-layer bidirectional \texttt{GRU} \cite{cho2014learning}, a single-layer bidirectional \texttt{LSTM} \cite{hochreiter1997long}, and a two-layer bidirectional \texttt{LSTM} with attention. The \emph{Transformer} encoders replace recurrence with self-attention \cite{vaswani2017attention} and learnable positional embeddings, ranging from 1 to 5 layers and 4 to 6 heads.

\iflncs\noindent\fi Finally, we assess transfer learning from pretrained language encoders. We consider frozen encoders with a light head trained on the pooled output: \texttt{BERT} \cite{devlin2019bert}, \texttt{RoBERTa} \cite{liu2019roberta}, \texttt{DeBERTa-v3} \cite{he2023debertav3}, \texttt{ELECTRA} \cite{clark2020electra}, \texttt{MPNet} \cite{song2020mpnet}, \texttt{ModernBERT} \cite{warner2024modernbert}, and \texttt{DistilBERT} \cite{sanh2019distilbert} on both text tasks, plus \texttt{MiniLM-L6} \cite{wang2020minilm} on \texttt{arXiv}. As with the vision encoders, freezing isolates pretrained representations from fine-tuning.
